무인 비행체 애드혹 네트워크(FANET)는 급격한 토폴로지 변화 속에서도 신뢰성 있고 저지연의 라우팅을 유지해야 하며, 자원이 제한된 UAV의 특성상 로컬 정보만을 활용한 분산형 실행이 필수적이다. 전역 정보나 에이전트 간 잦은 제어 메시지 교환에 의존하는 다중 에이전트 강화학습(MARL) 기법들은 우수한 경로를 학습할 수 있으나, 실제 분산 및 이동성 UAV 군집 환경에 배포하기에는 통신 오버헤드와 물리적 제약이라는 배포 격차(Deployment Gap)를 지닌다. 본 논문에서는 전역 그래프 정보를 지닌 교사(Teacher) 모델의 라우팅 정책을 1-hop 로컬 관측 정보만을 사용하는 학생(Student) 모델로 증류하는 정책 증류(Policy Distillation) 기법과, 링크 단절 및 버퍼 위험 상태에서만 예측적 라우팅 분기를 선택적으로 활성화하는 예측적 위험 스위칭(Predictive Risk Switching) 기반의 라우팅 프레임워크인 \method{}를 제안한다. 제안 기법은 평상시에는 가벼운 지리적 라우팅을 수행하여 지연 및 제어 오버헤드를 제어하며, 위험이 감지되는 예외 상황에서만 예측 모델을 구동하여 자원 효율성을 극대화한다. 총 14개 시나리오 및 5개 시드에 대한 시뮬레이션 평가 결과, \method{}의 전신인 \phaseTwelve{} 모델이 GPSR, 예측 지리적 라우팅, Evo-QGeo, IQMR Q($\lambda$), DRAMA 및 기존 분산 강화학습 모델 대비 제어 오버헤드 바이트를 현저히 적게 사용하면서도 가장 우수한 패킷 전달율(PDR) 및 전송 기한 준수율을 달성함을 확인하였다.
